\chapter{Ataxia}

\section*{Definition}
Ataxia is a neurological sign characterized by loss of coordination of voluntary movements, typically resulting from dysfunction of the cerebellum or its afferent/efferent pathways. It may affect gait, limb movements (dysmetria, intention tremor), speech (dysarthria), and eye movements (nystagmus), with onset acute (hours to days), subacute (days to weeks), or chronic (months to years).

\section*{What Happens in Your Body}
Ataxia arises when the cerebellum fails to coordinate the timing, force, and trajectory of movement. Cerebellar lesions disrupt the fine-tuning of descending motor commands and sensory feedback integration through the cortico-ponto-cerebellar and spinocerebellar pathways. Purkinje cell output through deep cerebellar nuclei modulates premotor cortex activity; damage causes decomposition of movement, dysmetria (overshoot/undershoot of target), and inability to perform rapid alternating movements (dysdiadochokinesia). Sensory (proprioceptive) ataxia results from disruption of dorsal column-medial lemniscal input to cerebellum.

\section*{Causes}
\begin{itemize}
 \item Acute ataxia: Cerebellar stroke (ischemic or hemorrhagic), Cerebellitis (post-infectious, especially varicella in children), Toxic exposures (ethanol intoxication, phenytoin, lithium, carbamazepine, benzodiazepines), Wernicke encephalopathy (thiamine deficiency), Multiple sclerosis relapse, Conversion disorder
 \item Subacute ataxia: Cerebellar tumor (medulloblastoma, astrocytoma, hemangioblastoma), Paraneoplastic syndromes (anti-Yo, anti-Hu, anti-Tr antibodies), Creutzfeldt-Jakob disease, HIV-associated cerebellar degeneration, Vitamin E deficiency, Hypothyroidism, Chronic alcohol-related cerebellar degeneration
 \item Chronic/progressive ataxia: Spinocerebellar ataxias (SCA1--SCA36, autosomal dominant), Friedreich ataxia (autosomal recessive, GAA repeat in FXN gene), Fragile X tremor/ataxia syndrome (FXTAS), Multiple system atrophy, Cerebellar type (MSA-C), Alexander disease
 \item Episodic ataxia: EA1 (KCNA1 mutation, with myokymia), EA2 (CACNA1A mutation, with nystagmus), Metabolic disorders (maple syrup urine disease, Hartnup disease, pyruvate dehydrogenase deficiency)
\end{itemize}

\section*{When to See a Doctor}

\begin{itemize}
 \item Acute onset of ataxia, especially with headache, vomiting, visual disturbance, or altered mental status (posterior circulation stroke/hemorrhage until proven otherwise)
 \item Ataxia with fever, nuchal rigidity, or altered consciousness (meningitis/encephalitis)
 \item Rapidly progressive ataxia over days to weeks (paraneoplastic or inflammatory)
 \item Ataxia following head trauma
 \item Child with acute gait ataxia (post-infectious cerebellar ataxia or posterior fossa tumor)
 \item Ataxia in the setting of known malignancy or immunocompromise
\end{itemize}

\section*{Related Symptoms}
\begin{itemize}
 \item Dysarthria (scanning speech)
 \item Nystagmus
 \item Dysmetria
 \item Intention tremor
 \item Dysdiadochokinesia
 \item Hypotonia
 \item Gait instability (wide-based gait)
 \item Vertigo
 \item Headache
 \item Cognitive impairment
\end{itemize}
\textbf{Important:} Acute ataxia you may need urgent brain imaging (CT or MRI) (CT or MRI brain) to exclude stroke, hemorrhage, or mass lesion; chronic progressive ataxia requires MRI, genetic testing, and screening for acquired causes.

\section*{Self-Care / Home Management}
\begin{itemize}
 \item Fall prevention: use of cane or walker, removal of tripping hazards, grab bars in bathroom, adequate lighting.
 \item Avoid alcohol entirely as it worsens cerebellar function.
 \item Physical and occupational therapy focusing on balance, gait, and coordination exercises.
 \item Discontinue non-prescribed medications that may cause or worsen ataxia (benzodiazepines, antiepileptics).
 \item Car driving should cease pending neurological evaluation and clearance.
\end{itemize}

\section*{How Your Doctor Will Evaluate You}
\begin{itemize}
 \item History determines temporal onset, progression pattern, alcohol/drug use, family history of ataxia, and associated symptoms (vertigo, headache, dysarthria).
 \item Neurological exam includes finger-to-nose, heel-to-shin, rapid alternating movements, gait assessment, Romberg test (distinguishes cerebellar from sensory ataxia), and eye movement examination.
 \item MRI brain with and without contrast is the imaging modality of choice for evaluating posterior fossa structures.
 \item Laboratory testing includes CBC, CMP, TSH, vitamin E, thiamine, toxicology screen, and paraneoplastic antibody panel when indicated.
 \item Genetic testing for SCA panel and Friedreich ataxia in those with positive family history or progressive ataxia with no acquired cause.
 \item Lumbar puncture for inflammatory/infectious etiologies when imaging unrevealing.
\end{itemize}

\section*{Treatment Options}
\begin{itemize}
 \item Acute ataxia: treat underlying cause (thiamine for Wernicke, IV steroids for MS flare-up, stroke management, discontinue offending drug).
 \item Immune-mediated ataxia: corticosteroids, IVIG, or plasma exchange for paraneoplastic or inflammatory cerebellitis.
 \item Tumor-associated ataxia: surgical resection, radiation, or chemotherapy depending on histology.
 \item Genetic ataxias: no disease-modifying therapy available; treatment is supportive (physical therapy, speech therapy, assistive devices).
 \item Symptomatic therapies: amantadine, buspirone, or riluzole may modestly improve coordination in some people; clonazepam for episodic ataxia type 2.
 \item Vitamin E supplementation for ataxia due to vitamin E deficiency (AVED or malabsorption).
\end{itemize}

\section*{References}
\begin{thebibliography}{99}
\bibitem{ataxia:ref1} Ashizawa T, Xia G. ``Ataxia.'' \textit{Continuum (Minneap Minn)}. 2016;22(4):1193-1218.
\bibitem{ataxia:ref2} Manto M, Gandini J, et al. ``Cerebellar ataxias: an update.'' \textit{Curr Opin Neurol}. 2020;33(4):506-513.
\bibitem{ataxia:ref3} Klockgether T, Mariotti C, et al. ``Spinocerebellar ataxia.'' \textit{Nat Rev Dis Primers}. 2019;5(1):24.
\bibitem{ataxia:ref4} Ropper AH, Samuels MA, et al. \textit{Adams and Victor\'s Principles of Neurology}, 11th ed. McGraw-Hill, 2019.
\bibitem{ataxia:ref5} UpToDate. ``Overview of cerebellar ataxia in adults.'' Wolters Kluwer, 2023.
\end{thebibliography}
